\documentclass[12pt,a4paper]{article}

\makeatletter
\def\input@path{{/home/maenwe/Master_CSM/modèle_cours_latex}}
\makeatother

\usepackage{prise_de_note}
\graphicspath{{images}}
\begin{document}
	\pagenumbering{gobble}
	\begin{titlepage}
		\centering
		
		% Logo
		\begin{figure}[h]
			\centering
			% Ajustez la largeur si besoin (ex: 0.4 ou 0.6)
			\includegraphics[width=0.5\textwidth]{/home/maenwe/Master_CSM/modèle_cours_latex/logo.jpg} 
		\end{figure}
		
		\vspace{2cm}
		
		% Filet horizontal du haut
		\rule{\linewidth}{0.5mm} \\[0.4cm]
		
		% Titre
		{\color{bleuTitre} \Huge \bfseries Phénomène de propagation} \\[0.4cm]
		
		% Filet horizontal du bas
		\rule{\linewidth}{0.5mm}
		
		\vspace{1.5cm}
		
		{\Large \textbf{Fiche de TD 2}}
		
		\vfill
		
		% Bloc Informations
		\begin{minipage}{0.45\textwidth}
			\begin{flushleft} \large
				\textbf{Enseignant :} \\
				M. Benjamin \textsc{BOUTIN}
			\end{flushleft}
		\end{minipage}
		\hfill
		\begin{minipage}{0.45\textwidth}
			\begin{flushright} \large
				\textbf{Étudiant :} \\
				Nathan \textsc{HASCOET}
			\end{flushright}
		\end{minipage}
		
		\vspace{2cm}
		
		% Date
		{\large \today}
		
	\end{titlepage}
	
	\newpage
	\pagenumbering{arabic}
	\setcounter{page}{1}

	\begin{sol}
		$ $\\
		 \begin{enumerate}
		 	\begin{minipage}{17cm}
		 		\item La matrice $A$ est de la forme : $A = \begin{bmatrix}
		 			0 & h_0\\
		 			g & 0
		 		\end{bmatrix}$. Les valeurs propres de $A$ sont les racines du polynôme $\det(XI-A) = X^2-gh_0$. Ainsi $\mathrm{sp}(A) = \{-c,c\}$.  Ainsi les deux vitesses caractéristiques sont $c$ et $-c$.
		 		
		 		\item On a : $\partial_t U = -i\omega r U$ et $\partial_x U = ikr U$. Ainsi l'équation caractéristique est $\det(-\omega I + kA) = 0$. Ce qui nous donne $\omega^2 = (kc)^2$. Ainsi $\omega_\pm = \pm kc$.
		 		
		 		\item Il n'est pas dispersif car $\frac{\omega_\pm}k$ ne dépend pas de $k$. De plus $\partial_tt \eta = -h0\partial_{xt} u = c^2 \partial_{xx} \eta$.
		 		Ainsi $\eta_{tt} - c^2 \nabla \eta = 0$. Donc $\eta$ vérifie l'équation des ondes.
		 		
		 		\item On pose $B = \begin{bmatrix}
		 			0 & 0\\
		 			0 & -\dfrac{h_0^2}3
		 		\end{bmatrix}$, et (BBM) se réécrit alors $\partial_t U + A\partial_x U+B\partial_{xxt} U = 0$. Ainsi pour $U = re^{i(kx-\omega t)}$, il suit : $-ir\omega U + ikr AU-ik^2 \omega r BU = 0$. Ainsi l'équation de dispersion est $\det(-\omega I + kA-k^2\omega B) = 0$. Ce qui revient à écrire :\\ $0 = \left|\begin{matrix}
		 		-\omega & kh_0\\
		 		kg & -\omega+\dfrac{h_0^2k^2\omega}3
		 		\end{matrix}\right| =\omega^2\left(1+\dfrac{h_0^2k^2}3\right) - k^2c^2$.
		 		
		 		\vspace{.5cm}
		 		
		 		\item $v_\varphi = \dfrac{\omega(k)}k = \dfrac{ck}{\sqrt{1+\dfrac{k^2h_0^2}3}} $, donc le modèle est dispersif.\\
		 		La vitesse de groupe est : $v_g(k) = \omega'(k) = \dfrac{c\sqrt{1+\dfrac{h_0^2k^2}3} -\dfrac{ch_0^2k^2}{3\sqrt{1+\dfrac{k^2h_0^2}3}}}{\left(1+\dfrac{k^2h_0^2}3\right)} = \dfrac{c}{\left(1+\dfrac{k^2h_0^2}3\right)^{3/2}}$
		 	\end{minipage}
		 \end{enumerate}
	\end{sol}

\end{document}